\scanpage{7}

\uline{\illegible{1}}\unsure{:|,} \unsure{soit} $\unsure{\mathcal{A}|\mathscr{A}|A}$ un anneau sur $X$, soient $M,N$ respectivement \unsure{les|les,} faisceaux $K^{0}$ et $K^{1}$ associés au champ additif des $\mathcal{A}$-Modules projectifs de v. f. \inserted{(\unsure{p|par} ex.),} Alors la \unsure{construction} de Mme Sinh nous fourni un champ de Picard épinglé par $M,N$, d'où une section \deleted{\unsure{du|de}} canonique du terme médian $\unsure{P|F}(M,N$ de $(*)$.

\uline{NB} Tout ce qui précède a les \unsure{fonctorialités} évidentes en $M,N,X$\ldots

\uline{Question} : Le triangle exact $(\unsure{T|J})$ \unsure{et} la suite exacte $(\unsure{*})$ sont ils \unsure{connu|connus} par les \unsure{compètences|compétences} (Quillen, Breen Illusie\ldots) ? \unsure{Connaissent-ils|Donnaissent-ils} des \unsure{variantes|variante} ``supérieures'' ? (Un principe ``géométrique'' pour en obtenir pourrait être via des $n$-champs de Picard non néc. stricts\ldots)

% Isolated point above the right of the Question's sequence label: intentional status uncertain; no textual function assigned.
% Isolated point above the gap between ils and connu: intentional status uncertain; no textual function assigned.
% Two curved impressions run along the lower left edge of the copied sheet; the surrounding straight frame is a copying edge.
% No manuscript page number is visible in the clipped opening.
% archive sheet number: 64
